\documentclass[11pt]{article}
\usepackage{mathblatt}
% vorspann.tex – Kompetenzblatt, zusammenbau v0.7 (2026-09-28).
% Ergänzt mathblatt.sty für Kompetenzblätter, ohne sie zu ändern
% (layout-befunde.md 1–34 vom 28.09.). Wird nach \usepackage{mathblatt}
% eingelesen.
\makeatletter
\RequirePackage{newunicodechar}
% Zeichen, die Latin Modern nicht hat (Befund 20): Text aus einer
% Ersatzschrift, Mathe als Befehl; ohne Ersatzschrift auch Text als Mathe.
\newif\ifkbersatz
\IfFontExistsTF{FreeSerif}{\newfontfamily\kbersatzschrift{FreeSerif}\kbersatztrue}{%
 \IfFontExistsTF{DejaVu Serif}{\newfontfamily\kbersatzschrift{DejaVu Serif}\kbersatztrue}{%
  \IfFontExistsTF{Cambria Math}{\newfontfamily\kbersatzschrift{Cambria Math}\kbersatztrue}{%
   \IfFontExistsTF{Segoe UI Symbol}{\newfontfamily\kbersatzschrift{Segoe UI Symbol}\kbersatztrue}{}}}}
\newcommand{\kbzeichen}[2]{\ifmmode #2\else\ifkbersatz{\kbersatzschrift #1}\else\ensuremath{#2}\fi\fi}
\newcommand{\kbzeichenhoch}[1]{\ifkbersatz\ifmmode\text{\kbersatzschrift #1}\else{\kbersatzschrift #1}\fi\else ?\fi}
% Kopf- und Fußzeile (Befund 1, 2): Kopf Thema · Blattart, Fuß Kennung und Seite
\newcommand{\kbkopfzeile}[2]{\pagestyle{fancy}\fancyhf{}\renewcommand{\headrulewidth}{0pt}%
  \fancyhead[L]{\footnotesize\color{mbgrau}#1}%
  \fancyfoot[L]{\footnotesize\color{mbgrau}#2}%
  \fancyfoot[R]{\footnotesize\color{mbgrau}Seite \thepage}}
% Flattersatz überall (Befund 25)
\raggedright
% Titel des Blatts: Ich-kann-Satz, darunter Zweigzeile
\newcommand{\kbtitel}[2]{\par\noindent{\Large\bfseries\raggedright #1\par}%
  \ifblank{#2}{}{\par\smallskip\noindent{\small #2}\par}\medskip}
% Abschnitt: „Das kennst du schon“, „Zum Merken“
\newcommand{\kbabschnitt}[1]{\par\addvspace{12pt}\Needspace*{5\baselineskip}%
  \noindent{\bfseries #1}\par\nobreak\vspace{1pt}\noindent{\color{mbgrau}\rule{\linewidth}{0.4pt}}\par\nobreak}
% Hauptnummer: Titel hängt am ersten Block; jeder Block (Teilaufgabe) bricht
% nicht, passt er nicht mehr, rückt er ganz auf die nächste Seite (Befund 21).
\newif\ifkbtitel
\newsavebox{\kbbox}
\newlength{\kbhoehe}
\newlength{\kbnummerabstand}\setlength{\kbnummerabstand}{10pt}
\newlength{\kbteilabstand}\setlength{\kbteilabstand}{6pt}
\newenvironment{kbaufgabe}[1]{\stepcounter{aufgabe}\setcounter{teil}{0}%
  \gdef\kbtiteltext{\noindent\hangindent1.6em\hangafter1\makebox[1.6em][l]{\textbf{\theaufgabe.}}#1\par}%
  \global\kbtiteltrue\par\addvspace{\kbnummerabstand}}{\par}
\newenvironment{kbblock}{\par\addvspace{\kbteilabstand}\begin{lrbox}{\kbbox}%
  \begin{minipage}[t]{\linewidth}\raggedright\parskip\z@
  \ifkbtitel\kbtiteltext\vspace{2pt}\global\kbtitelfalse\fi
  \list{}{\leftmargin1.6em\labelwidth1.4em\labelsep0.2em\topsep\z@\partopsep\z@
    \parsep\z@\itemsep\z@\rightmargin\z@}\raggedright}%
  {\endlist\end{minipage}\end{lrbox}%
  \setlength{\kbhoehe}{\dimexpr\ht\kbbox+\dp\kbbox\relax}%
  \Needspace*{\kbhoehe}\noindent\usebox{\kbbox}\par}
% Paarsatz (Platz nutzen, Befund 27): zwei Hauptnummern oder zwei
% Teilaufgaben mit Zeichenaufgabe nebeneinander, als ein Block
\newcommand{\kbnummer}[1]{\stepcounter{aufgabe}\setcounter{teil}{0}%
  \noindent\hangindent1.6em\hangafter1\makebox[1.6em][l]{\textbf{\theaufgabe.}}#1\par\vspace{2pt}}
\newcommand{\kbhalb}[1]{\list{}{\leftmargin1.6em\labelwidth1.4em\labelsep0.2em\topsep\z@
  \partopsep\z@\parsep\z@\itemsep\z@}\raggedright #1\endlist}
\newcommand{\kbpaar}[2]{\par\addvspace{\kbnummerabstand}\begin{lrbox}{\kbbox}%
  \begin{minipage}[t]{0.485\linewidth}\raggedright\parskip\z@ #1\end{minipage}\hfill
  \begin{minipage}[t]{0.485\linewidth}\raggedright\parskip\z@ #2\end{minipage}\end{lrbox}%
  \setlength{\kbhoehe}{\dimexpr\ht\kbbox+\dp\kbbox\relax}\Needspace*{\kbhoehe}\noindent\usebox{\kbbox}\par}
\newcommand{\kbteilpaar}[2]{\item[]\noindent
  \begin{minipage}[t]{0.485\linewidth}\kbhalb{#1}\end{minipage}\hfill
  \begin{minipage}[t]{0.485\linewidth}\kbhalb{#2}\end{minipage}\par}
% Anweisung der Hauptnummer, eigene Zeile über der Teilaufgabe (Befund 5)
\newcommand{\kbanweisung}[1]{\item[]#1\par\vspace{2pt}}
% Kopfzeile einer Teilaufgabe: Buchstabe, Auftakt (halbfett oder Kontext),
% Prüfkennung klein rechts (Befund 3, 12, Beschluss c)
\newlength{\kbkennbreite}\setlength{\kbkennbreite}{1.9cm}
\newcommand{\kbkennung}[1]{{\footnotesize #1}}
\newcommand{\kbteil}[3]{\item[#1]\ifblank{#3}%
  {\parbox[t]{\linewidth}{\raggedright #2}}%
  {\parbox[t]{\dimexpr\linewidth-\kbkennbreite\relax}{\raggedright #2}\hfill
   \parbox[t]{\kbkennbreite}{\raggedleft\kbkennung{#3}}}\par}
% Frage in eigener Zeile (Befund 4), ein Satz je Zeile (Befund 10)
\newcommand{\kbfrage}[1]{\par\vspace{2pt}#1\par}
% Antwortfeld in eigener Zeile darunter, linksbündig (Befund 4, 8)
\newcommand{\kbantwort}[1]{\par\vspace{4pt}\noindent#1\par}
% Zwei Spalten: links Grafik/Tabelle/Aufgabe, rechts Antwort oder Rechenraum
% (Befund 27); oben bündig. Beginnt einen eigenen Listenpunkt ohne Marke.
\newcommand{\kbzweispaltig}[3][0.48]{\item[]\noindent
  \begin{minipage}[t]{#1\linewidth}\raggedright #2\end{minipage}\hfill
  \begin{minipage}[t]{\dimexpr0.98\linewidth-#1\linewidth\relax}\raggedright #3\end{minipage}\par}
% Linke Spalte mit Teilaufgabe (Buchstabe hängt links heraus wie in der Liste)
\newcommand{\kbspalte}[1]{\list{}{\leftmargin\z@\labelwidth1.4em\labelsep0.2em\topsep\z@
  \partopsep\z@\parsep\z@\itemsep\z@}\raggedright #1\endlist}
% Antwortfelder: 2,2 cm, in Punkten 1,2 cm; ohne Umbruch davor
\newcommand{\kbfeld}[1][]{\mbox{\underline{\hspace{2.2cm}}\ifblank{#1}{}{\,#1}}}
\newcommand{\kbfeldk}[1][]{\mbox{\underline{\hspace{1.2cm}}\ifblank{#1}{}{\,#1}}}
% Grafik oder Tabelle unter dem Text, linksbündig (Befund 8, 28)
\newcommand{\kbdarunter}[1]{\par\vspace{4pt}\noindent#1\par}
% Ankreuzen: je Option eine Zeile, linksbündig (Befund 6, 7)
\newcommand{\kbkreuz}[1]{\par\vspace{2pt}\noindent$\square$\ #1\par}
% Bearbeitungsraum (Befund 9): Schreibzeilen wie mathblatt (9 mm)
\newcommand{\kbraum}[1]{\schreibzeilen{#1}}
% Wertetabelle mit Köpfen im Textmodus (Befund 26): wie \wertetabelle der
% Vorlage, aber die Köpfe setzt der Aufruf selbst ($x$, „Zeit in h“).
\NewDocumentCommand{\kbwertetabelle}{o m m m}{%
  \def\mbkopf{}\def\mbwerte{}\def\mbleer{}\setcounter{mbn}{0}%
  \foreach \v in {#4}{\xappto\mbkopf{& $\v$}\xappto\mbleer{& }\stepcounter{mbn}\xdef\mbnn{\arabic{mbn}}}%
  \IfNoValueTF{#1}{}{\foreach \v in {#1}{\ifdefempty{\v}{\xappto\mbwerte{& }}{\xappto\mbwerte{& $\v$}}}}%
  \begingroup\renewcommand{\arraystretch}{1.3}%
  \edef\mbpre{|>{\noexpand\columncolor{mbkasten}}l|*{\mbnn}{>{\noexpand\centering\noexpand\arraybackslash}p{\noexpand\mbzell}|}}%
  \expandafter\mbtabstart\expandafter{\mbpre}\hline
  #2 \mbkopf \\ \hline
  \IfNoValueTF{#1}{\rule{0pt}{\mbzeile}#3 \mbleer \\ \hline}{#3 \mbwerte \\ \hline}
  \end{tabular}\endgroup}
% Merkkasten am Ende (Aufbau des Kompetenzblatts)
\newcommand{\kbmerk}[2]{\par\addvspace{14pt}\begin{lrbox}{\kbbox}\begin{minipage}[t]{\linewidth}%
  \noindent{\bfseries #1}\par\vspace{1pt}\noindent{\color{mbgrau}\rule{\linewidth}{0.4pt}}\par\vspace{4pt}%
  \uebersichtskasten{\raggedright #2}\end{minipage}\end{lrbox}%
  \setlength{\kbhoehe}{\dimexpr\ht\kbbox+\dp\kbbox\relax}\Needspace*{\kbhoehe}\noindent\usebox{\kbbox}\par}
% Lösungsblatt: Nummer und Buchstabe halbfett, Lösung daneben
\newcommand{\kbloesung}[2]{\par\addvspace{3pt}\noindent\hangindent2.8em\hangafter1\makebox[2.8em][l]{\textbf{#1}}#2\par}
\makeatother

\begin{document}
\kbkopfzeile{Lineare Funktionen · Kompetenzblatt}{LIN-K1}
% Kompetenzblatt LIN-K1: lineare-funktionen e2 „Graph zeichnen“ (zusammenbau v0.7, Niveau FOR)
\kbtitel{Ich kann den Graphen einer linearen Funktion zeichnen.}{ab Kl. 8 · P10 ×5}
\setcounter{aufgabe}{0}

\kbabschnitt{Das kennst du schon}

% zone: lineare-funktionen-zone-f4-v1
\begin{kbaufgabe}{Ich kann einen Bruch mit einer ganzen Zahl malnehmen.}
\begin{kbblock}
\kbanweisung{Rechne.}
\kbteil{}{{\bfseries\boldmath $\frac{1}{2} \cdot 8$}}{}
\kbantwort{\kbfeld}
\end{kbblock}
\end{kbaufgabe}

% zone: lineare-funktionen-zone-f1-v1
\begin{kbaufgabe}{Ich kann Koordinaten ablesen und eintragen.}
\begin{kbblock}
\kbteil{}{{\bfseries\boldmath Punkt A}}{}
\kbfrage{Welche Koordinaten? \\ Lies ab.}
\kbzweispaltig[0.47]{\begin{ksys}[xmin=-5,xmax=5,ymin=-5,ymax=5,ablesen] \punkt{4}{3}{A} \end{ksys}}{\kbantwort{A( \kbfeldk | \kbfeldk )}}
\end{kbblock}
\end{kbaufgabe}

% zone: lineare-funktionen-zone-f3-v1
\begin{kbaufgabe}{Ich kann mit negativen Zahlen malnehmen und teilen.}
\begin{kbblock}
\kbanweisung{Rechne.}
\kbteil{}{{\bfseries\boldmath $(-3) \cdot 4$}}{}
\kbantwort{\kbfeld}
\end{kbblock}
\end{kbaufgabe}

\kbabschnitt{Schritt für Schritt}

% leiter: lineare-funktionen-e2-k4-s1-v1
\begin{kbaufgabe}{Ich kann zu einer Funktion eine Wertetabelle ausfüllen.}
\begin{kbblock}
\kbteil{}{{\bfseries\boldmath Wertetabelle}}{}
\kbfrage{Fülle die Tabelle zu $f(x) = 3x - 2$ aus.}
\kbdarunter{\wertetabelle{x}{f(x)}{-1,0,1,2}}
\end{kbblock}
\end{kbaufgabe}

% leiter: lineare-funktionen-e2-k4-s2-v1, lineare-funktionen-e2-k4-s3-v1 (Paar)
\kbpaar{
\kbnummer{Ich kann Wertepaare als Punkte eintragen.}\kbhalb{\kbteil{}{{\bfseries\boldmath Punkte eintragen}}{}
\kbfrage{Trage die Wertepaare $(-2|-2)$, $(-1|1)$, $(0|4)$, $(1|7)$ der Funktion $f(x) = 3x + 4$ ins Koordinatensystem ein.}
\kbdarunter{\begin{ksys}[xmin=-4,xmax=4,ymin=-4,ymax=8,karo=0.56] \end{ksys}}}
}{
\kbnummer{Ich kann zwei Punkte berechnen und die Gerade durch sie zeichnen.}\kbhalb{\kbteil{}{{\bfseries\boldmath Gerade durch Punkte}}{}
\kbfrage{Berechne zwei Punkte von $f(x) = 4x + 2$, trage sie ein und zeichne die Gerade durch sie.}
\kbdarunter{\begin{ksys}[xmin=-4,xmax=4,ymin=-4,ymax=7,karo=0.61] \end{ksys}}}
}

% leiter: lineare-funktionen-e2-k4-s4-v1, lineare-funktionen-e2-k4-s5-v1 (Paar)
\kbpaar{
\kbnummer{Ich kann die Gerade mit n und dem Steigungsdreieck zeichnen.}\kbhalb{\kbteil{}{{\bfseries\boldmath Mit n und Steigungsdreieck}}{}
\kbfrage{Zeichne die Gerade zu $f(x) = 3x + 5$: Markiere n auf der y-Achse, gehe 1 nach rechts und m nach oben.}
\kbdarunter{\begin{ksys}[xmin=-4,xmax=4,ymin=-4,ymax=9,karo=0.52] \end{ksys}}}
}{
\kbnummer{Ich kann eine fallende Gerade zeichnen, wenn m negativ ist.}\kbhalb{\kbteil{}{{\bfseries\boldmath Mit n und Steigungsdreieck}}{}
\kbfrage{Zeichne die Gerade zu $f(x) = -3x + 4$.}
\kbdarunter{\begin{ksys}[xmin=-4,xmax=4,ymin=-4,ymax=5,karo=0.74] \end{ksys}}}
}

% leiter: lineare-funktionen-e2-k4-s6-v1, lineare-funktionen-e2-k4-s7-v1 (Paar)
\kbpaar{
\kbnummer{Ich kann die Gerade zeichnen, wenn m ein Bruch ist.}\kbhalb{\kbteil{}{{\bfseries\boldmath Mit n und Steigungsdreieck}}{}
\kbfrage{Zeichne die Gerade zu $f(x) = \frac{1}{2}x + 3$.}
\kbdarunter{\begin{ksys}[xmin=-4,xmax=4,ymin=-4,ymax=5,karo=0.74] \end{ksys}}}
}{
\kbnummer{Ich kann die Gerade zeichnen, wenn n negativ ist.}\kbhalb{\kbteil{}{{\bfseries\boldmath Mit n und Steigungsdreieck}}{}
\kbfrage{Zeichne die Gerade zu $f(x) = 2x - 4$.}
\kbdarunter{\begin{ksys}[xmin=-4,xmax=4,ymin=-5,ymax=4,karo=0.74] \end{ksys}}}
}

% pruefung: lineare-funktionen-e2-k4-s8-v1, lineare-funktionen-e2-k4-s8-v3, lineare-funktionen-e2-k4-s8-v5, lineare-funktionen-e2-k4-s8-v7
\begin{kbaufgabe}{Ich kann das auch in Aufgaben aus der Prüfung.}
\begin{kbblock}
\kbteilpaar{
\kbteil{a)}{{\bfseries\boldmath Graph zeichnen}}{P10 ’26 F}
\kbfrage{Zeichne den Graphen der Funktion f mit $f(x) = -4x + 3$ in das Koordinatensystem.}
\kbdarunter{\begin{ksys}[xmin=-4,xmax=4,ymin=-4,ymax=4] \end{ksys}}
}{
\kbteil{b)}{{\bfseries\boldmath Graph zeichnen}}{P10 ’24}
\kbfrage{Zeichne den Graphen der Funktion f mit $f(x) = 5x - 3$ in das Koordinatensystem.}
\kbdarunter{\begin{ksys}[xmin=-4,xmax=4,ymin=-4,ymax=4] \end{ksys}}
}
\end{kbblock}
\begin{kbblock}
\kbteilpaar{
\kbteil{c)}{{\bfseries\boldmath Graph zeichnen}}{P10 ’22}
\kbfrage{Zeichne den Graphen der Funktion f mit $f(x) = -3x + 5$ in das Koordinatensystem.}
\kbdarunter{\begin{ksys}[xmin=-4,xmax=4,ymin=-4,ymax=6,karo=0.67] \end{ksys}}
}{
\kbteil{d)}{{\bfseries\boldmath Graph zeichnen}}{P10 ’21}
\kbfrage{Zeichne den Graphen der Funktion f mit $f(x) = 2x + 5$ in das Koordinatensystem.}
\kbdarunter{\begin{ksys}[xmin=-4,xmax=4,ymin=-4,ymax=8,karo=0.56] \end{ksys}}
}
\end{kbblock}
\end{kbaufgabe}

\kbmerk{Zum Merken}{Lineare Funktion: f(x) = m · x + n \\ m = Steigung: 1 nach rechts, m nach oben (m < 0: nach unten). \quad f(x) = 2x + 1: 1 nach rechts, 2 nach oben \\ n = y-Achsenabschnitt: die Gerade schneidet die y-Achse bei (0 | n). \quad f(x) = 2x + 1: bei (0 | 1)}
\end{document}
